\documentclass[12pt,a4paper]{article}
\usepackage[fontset=none]{ctex}
\setCJKmainfont{SimSun}

\usepackage[margin=1in]{geometry}
\usepackage{amsmath,amssymb,bm}
\usepackage{booktabs}
\usepackage{graphicx}
\usepackage{hyperref}
\usepackage{fancyhdr}

\hypersetup{colorlinks,linkcolor=blue,citecolor=blue,urlcolor=blue}
\setlength{\headheight}{15pt}
\setlength{\emergencystretch}{3em}

\pagestyle{fancy}
\fancyhf{}
\fancyhead[L]{二维浅水方程 4D-Var 离散伴随}
\fancyhead[R]{\thepage}

\title{二维浅水方程 4D-Var 离散伴随说明\\
\large 原版程序的历史说明（见 \texttt{baseline/original.c}）}
\author{杨鹏亮}
\date{\today}

\begin{document}

\maketitle

\noindent\textbf{版本说明（2026-10-16）：}本文保留为原版程序的推导和实验记录，
不描述现行实现。现行代码使用检查点重算伴随、相关背景协方差和白化控制变量，
以 SciPy L-BFGS-B 优化；配置、验证和新实验结果见 \texttt{README.md}、
\texttt{RESULTS.md} 与 \texttt{var\_output/summary.json}。


\begin{abstract}
本文只记录当前程序 \texttt{swe2d\_4dvar.c} 实际实现的离散伴随。
前向模式是二维浅水方程的有限体积离散，界面通量为 Rusanov 通量，
时间推进为两次 Euler 子步组成的 RK2。4D-Var 控制变量是初始全状态，
观测只抽取若干网格点上的水深。梯度由完整前向轨道、逐步反向的
离散伴随和背景项共同组成。本文不再推导连续伴随 PDE，因为代码没有
使用连续伴随后再离散的路线。
\end{abstract}

\tableofcontents

\section{状态、网格和代价函数}

单元状态为守恒变量
\[
\mathbf{q}_{i,j} =
\begin{pmatrix}
h_{i,j}\\
(hu)_{i,j}\\
(hv)_{i,j}
\end{pmatrix}.
\]
代码中一个完整场用三个数组 \texttt{h}, \texttt{hu}, \texttt{hv} 保存，
单元索引为
\[
k=i+N_x j.
\]

当前 4D-Var 只优化初始场 $\mathbf{Q}^0\in\mathbb{R}^{3N}$。背景场记为
$\mathbf{Q}_b$，观测位于每个同化窗口末端。代价函数为
\begin{equation}
J(\mathbf{Q}^0)
=
\frac12(\mathbf{Q}^0-\mathbf{Q}_b)^T B^{-1}(\mathbf{Q}^0-\mathbf{Q}_b)
+\frac12\sum_{c=1}^{N_c}
\left(H\mathbf{Q}^{n_c}-\mathbf{y}_c\right)^T
R^{-1}
\left(H\mathbf{Q}^{n_c}-\mathbf{y}_c\right).
\label{eq:cost}
\end{equation}
这里 $N=N_xN_y$，$n_c=c\,\mathtt{VAR\_STEPS\_WINDOW}$。
$H$ 是观测算子，从全状态中抽取观测单元的水深 $h$ 分量：
\[
H\mathbf{Q}^{n_c}
=
\left(h_{o_1}^{n_c},\ldots,h_{o_{N_{\rm obs}}}^{n_c}\right)^T,
\]
其中 $o_p$ 是第 $p$ 个观测单元索引。

矩阵 $B$ 是背景误差协方差矩阵；当前代码采用逐格点、逐变量独立的对角形式
\[
B=\operatorname{diag}\left(
\sigma_h^2 I_N,\,
\sigma_m^2 I_N,\,
\sigma_m^2 I_N
\right).
\]
矩阵 $R$ 是观测误差协方差矩阵；当前代码假设所有水深观测误差相互独立且方差相同，
\[
R=\sigma_o^2 I_{N_{\rm obs}}.
\]
因此式 \eqref{eq:cost} 展开后正是代码中
\texttt{evaluate\_cost\_gradient()} 计算的逐分量加权平方和。

\section{前向离散模式}

\subsection{物理通量}

令 $m=hu$，$n=hv$。物理通量为
\begin{equation}
\mathbf{f}(\mathbf{q}) =
\begin{pmatrix}
m\\
m^2/h+\frac12 g h^2\\
mn/h
\end{pmatrix},
\qquad
\mathbf{g}(\mathbf{q}) =
\begin{pmatrix}
n\\
mn/h\\
n^2/h+\frac12 g h^2
\end{pmatrix}.
\label{eq:physical-flux}
\end{equation}
代码在分母中使用 \texttt{safe\_h()}，即 $h$ 小于 \texttt{HMIN} 时用
\texttt{HMIN} 代替。

\subsection{Rusanov 通量}

$x$ 方向界面通量为
\begin{equation}
\widehat{\mathbf{f}}(\mathbf{q}_L,\mathbf{q}_R)
=
\frac12\left[\mathbf{f}(\mathbf{q}_L)+\mathbf{f}(\mathbf{q}_R)\right]
-\frac12 s(\mathbf{q}_R-\mathbf{q}_L),
\label{eq:rusanov-x}
\end{equation}
其中
\[
s=\max(a_L,a_R),
\qquad
a(\mathbf{q})=|u|+\sqrt{g h},
\qquad
u=m/h.
\]
$y$ 方向完全类似，只是速度用 $v=n/h$：
\[
a_y(\mathbf{q})=|v|+\sqrt{g h}.
\]
这些式子对应 \texttt{rusanov\_x()} 和 \texttt{rusanov\_y()}。

\subsection{Euler 子步}

\texttt{euler\_update()} 实现
\begin{equation}
\mathbf{q}_{i,j}^{+}
=
\mathbf{q}_{i,j}
-\frac{\Delta t}{\Delta x}
\left(
\widehat{\mathbf{f}}_{i+1/2,j}
-\widehat{\mathbf{f}}_{i-1/2,j}
\right)
-\frac{\Delta t}{\Delta y}
\left(
\widehat{\mathbf{g}}_{i,j+1/2}
-\widehat{\mathbf{g}}_{i,j-1/2}
\right).
\label{eq:euler}
\end{equation}
越界邻居由 \texttt{get\_state\_bc()} 钳制到最近的边界单元，因此边界是
离散零梯度延拓。

\subsection{RK2 步}

\texttt{step\_rk2\_store()} 使用
\begin{align}
\mathbf{Q}^{(1)} &= E(\mathbf{Q}^n),\\
\mathbf{Q}^{(2)} &= E(\mathbf{Q}^{(1)}),\\
\mathbf{Q}^{n+1} &= \frac12\mathbf{Q}^{n}+\frac12\mathbf{Q}^{(2)}.
\label{eq:rk2}
\end{align}
其中 $E$ 是式 \eqref{eq:euler} 的 Euler 子步。代码保存
\texttt{traj[n]} 和 \texttt{stage1[n]}，分别对应 $\mathbf{Q}^n$ 和
$\mathbf{Q}^{(1)}$，供离散伴随反向使用。

\section{界面通量的 Jacobian}

离散伴随需要 Rusanov 通量对两侧状态的导数。物理通量 Jacobian 为
\begin{equation}
A(\mathbf{q})=\frac{\partial\mathbf{f}}{\partial\mathbf{q}}
=
\begin{pmatrix}
0 & 1 & 0\\
gh-m^2/h^2 & 2m/h & 0\\
-mn/h^2 & n/h & m/h
\end{pmatrix},
\label{eq:Ax}
\end{equation}
\begin{equation}
B(\mathbf{q})=\frac{\partial\mathbf{g}}{\partial\mathbf{q}}
=
\begin{pmatrix}
0 & 0 & 1\\
-mn/h^2 & n/h & m/h\\
gh-n^2/h^2 & 0 & 2n/h
\end{pmatrix}.
\label{eq:Ay}
\end{equation}
代码中的 \texttt{flux\_jacobian\_x()} 和 \texttt{flux\_jacobian\_y()}
还包含 \texttt{safe\_h()} 分支：当原始 $h\le\mathtt{HMIN}$ 时，关于
$h$ 的导数被置零，以匹配前向分支。

对 $x$ 方向 Rusanov 通量，定义
\[
\Delta\mathbf{q}=\mathbf{q}_R-\mathbf{q}_L.
\]
若 $a_L\ge a_R$，则 $s=a_L$，$\nabla_{\mathbf{q}_L}s=\nabla a_L$，
$\nabla_{\mathbf{q}_R}s=0$；否则相反。于是
\begin{align}
\widehat{F}_L
&=
\frac12 A(\mathbf{q}_L)+\frac12sI
-\frac12\Delta\mathbf{q}\,(\nabla_{\mathbf{q}_L}s)^{T},
\label{eq:FL}\\
\widehat{F}_R
&=
\frac12 A(\mathbf{q}_R)-\frac12sI
-\frac12\Delta\mathbf{q}\,(\nabla_{\mathbf{q}_R}s)^{T}.
\label{eq:FR}
\end{align}
这正是 \texttt{rusanov\_x\_jacobians()} 的计算。$y$ 方向的
\texttt{rusanov\_y\_jacobians()} 用同样公式，只把左右状态换成
bottom/top 状态，并用 $B(\mathbf{q})$ 和 $a_y$。

速度导数在代码中为
\begin{equation}
\nabla a_x =
\begin{pmatrix}
-\operatorname{sgn}(u)m/h^2+\frac12 g/\sqrt{gh}\\
\operatorname{sgn}(u)/h\\
0
\end{pmatrix},
\qquad
\nabla a_y =
\begin{pmatrix}
-\operatorname{sgn}(v)n/h^2+\frac12 g/\sqrt{gh}\\
0\\
\operatorname{sgn}(v)/h
\end{pmatrix}.
\label{eq:wavespeed-grad}
\end{equation}
在 $u=0$ 或 $v=0$ 时，\texttt{signum()} 返回零。

\section{Euler 子步的离散伴随}

设某个输出单元 $(i,j)$ 的伴随为
\[
\bm{\lambda}_{i,j}^{+}
=
\begin{pmatrix}
\lambda_h^+\\
\lambda_{hu}^+\\
\lambda_{hv}^+
\end{pmatrix}.
\]
\texttt{euler\_adjoint()} 对每个输出单元做反向累加。首先，前向 Euler
子步中有恒等项 $\mathbf{q}_{i,j}^{+}\leftarrow\mathbf{q}_{i,j}$，因此
\[
\bm{\lambda}_{i,j}\mathrel{+}= \bm{\lambda}_{i,j}^{+}.
\]

然后四个界面按前向式 \eqref{eq:euler} 中的符号回传：
\begin{align}
\text{右界面:}\quad
\bm{\lambda}_{i,j}   &\mathrel{+}= -\frac{\Delta t}{\Delta x}\widehat{F}_L^T\bm{\lambda}_{i,j}^{+},
&
\bm{\lambda}_{i+1,j} &\mathrel{+}= -\frac{\Delta t}{\Delta x}\widehat{F}_R^T\bm{\lambda}_{i,j}^{+},
\label{eq:adj-right}\\
\text{左界面:}\quad
\bm{\lambda}_{i-1,j} &\mathrel{+}=  \frac{\Delta t}{\Delta x}\widehat{F}_L^T\bm{\lambda}_{i,j}^{+},
&
\bm{\lambda}_{i,j}   &\mathrel{+}=  \frac{\Delta t}{\Delta x}\widehat{F}_R^T\bm{\lambda}_{i,j}^{+},
\label{eq:adj-left}\\
\text{上界面:}\quad
\bm{\lambda}_{i,j}   &\mathrel{+}= -\frac{\Delta t}{\Delta y}\widehat{G}_B^T\bm{\lambda}_{i,j}^{+},
&
\bm{\lambda}_{i,j+1} &\mathrel{+}= -\frac{\Delta t}{\Delta y}\widehat{G}_T^T\bm{\lambda}_{i,j}^{+},
\label{eq:adj-top}\\
\text{下界面:}\quad
\bm{\lambda}_{i,j-1} &\mathrel{+}=  \frac{\Delta t}{\Delta y}\widehat{G}_B^T\bm{\lambda}_{i,j}^{+},
&
\bm{\lambda}_{i,j}   &\mathrel{+}=  \frac{\Delta t}{\Delta y}\widehat{G}_T^T\bm{\lambda}_{i,j}^{+}.
\label{eq:adj-bottom}
\end{align}
矩阵转置乘向量由 \texttt{add\_matT\_vec\_to\_cell()} 完成。

边界处使用和前向相同的钳制索引。例如 $i=0$ 时，左邻居仍然是
$i=0$；因此边界单元会收到来自自身和钳制 ghost 路径的累加贡献。
这不是连续边界条件的近似，而是前向离散索引映射的精确转置。

\section{RK2 步的离散伴随}

由式 \eqref{eq:rk2} 可得反向顺序：
\begin{align}
\bm{\Lambda}^{(2)} &= \frac12\bm{\Lambda}^{n+1},\\
\bm{\Lambda}^{n}   &= \frac12\bm{\Lambda}^{n+1},\\
\bm{\Lambda}^{(1)} &= E'(\mathbf{Q}^{(1)})^T\bm{\Lambda}^{(2)},\\
\bm{\Lambda}^{n}   &\mathrel{+}= E'(\mathbf{Q}^{n})^T\bm{\Lambda}^{(1)}.
\label{eq:rk2-adj}
\end{align}
这对应 \texttt{rk2\_adjoint\_step()}：
\begin{itemize}
\item \texttt{lam\_stage2} 接收 $\frac12\bm{\Lambda}^{n+1}$；
\item \texttt{lam\_n} 先接收 $\frac12\bm{\Lambda}^{n+1}$；
\item 在 \texttt{Q1} 上调用一次 \texttt{euler\_adjoint()}；
\item 在 \texttt{Qn} 上再调用一次 \texttt{euler\_adjoint()}；
\item 最后把第二次反传结果累加到 \texttt{lam\_n}。
\end{itemize}

\section{观测项和全窗口梯度}

观测算子只抽取水深：
\[
H\mathbf{Q}^{n_c}
=
\left(
h_{o_1}^{n_c},\ldots,h_{o_{N_{\rm obs}}}^{n_c}
\right)^T.
\]
因此在观测时刻只向 $h$ 分量注入
\begin{equation}
\lambda_h(o_p)\mathrel{+}=
\frac{h_{o_p}^{n_c}-y_{c,p}}{\sigma_o^2}.
\label{eq:obs-adj}
\end{equation}
这对应 \texttt{obs\_gradient\_add()}。

\texttt{evaluate\_cost\_gradient()} 的顺序是：
\begin{enumerate}
\item 从当前初值 \texttt{x0} 前向积分，保存 \texttt{traj} 和 \texttt{stage1}；
\item 计算背景项，并把背景项梯度写入 \texttt{grad}；
\item 遍历所有观测时刻，累计观测代价；
\item 从末时刻反向到初始时刻；若当前时间层是观测时刻，先执行式
      \eqref{eq:obs-adj}，再反传一个 RK2 步；
\item 把回到 $t=0$ 的伴随 \texttt{lam\_curr} 加到背景项梯度上。
\end{enumerate}
最终梯度为
\begin{equation}
\nabla_{\mathbf{Q}^0}J
=
B^{-1}(\mathbf{Q}^0-\mathbf{Q}_b)+\bm{\Lambda}^{0},
\label{eq:grad}
\end{equation}
其中 $B$ 是对角矩阵，每个单元对应
\[
\operatorname{diag}(\sigma_h^2,\sigma_m^2,\sigma_m^2).
\]


\section{结论}

当前程序实现的是离散伴随，不是连续伴随 PDE 的再离散。保持一致的关键是：
\begin{enumerate}
\item 使用和前向完全相同的 Rusanov 分支、边界钳制和固定时间步长；
\item 对 Euler 子步逐界面累加 Jacobian 转置贡献；
\item 对 RK2 严格按组合映射反向传播；
\item 在观测时刻只向被观测水深分量注入 $R^{-1}(H\mathbf{Q}-\mathbf{y})$；
\item 最后把初始伴随加到背景项梯度上。
\end{enumerate}
这些步骤与 \texttt{swe2d\_4dvar.c} 中的函数调用顺序一一对应。

\end{document}
